\scanpage{419}
\mspage{158}
à l'aide des deux flèches canoniques
\diagram{\begin{tikzcd}[ampersand replacement=\&]
\mathcal C_{\mathbb N^{\unsure{o|0}}}\arrow[r]\arrow[d] \& \underline{\mathrm{Ch}}_{\infty}(X)\\
\underline{\mathrm{Ch}}(X) \& {,}
\end{tikzcd}}
on sera obligé de prouver aussi les résultats qui correspondent à 1) 2) pour $\mathcal C_{\mathbb N^o}$, i.e.

% label at left of the diagram
\diagram{\begin{tikzcd}[ampersand replacement=\&,row sep=small]
{1')}\& X\arrow[r] \& \mathcal C_{\mathbb N^o}\\
\& \mathcal C_{t_0}=\unsure{C|\mathcal C}_0\arrow[u,phantom,"\simeq"{description,rotate=90}] \&
\end{tikzcd}}
(l'inclusion canonique de $C_0$, fibre de $\mathcal C_{\mathbb N^o}$ (fibré \unsure{en|sur} $\mathbb N^o$) en $0\in\mathbb N^o$) est une Hot-équivalence, et
\begingroup\makeatletter\tagsleft@true\makeatother
\[
\mathcal C_{\mathbb N^o}\longrightarrow X\times X\qquad\text{est propre et lisse.}\tag*{$2')$}
\]
\endgroup

Moyennant ces \deleted{\unsure{rés}} quatre résultats techniques, que je vais prouver tantôt (et $1'$ et $2'$ ne font aucune surprise pour moi), on gagne. Considérons en effet le diagramme
\diagram{\begin{tikzcd}[ampersand replacement=\&]
\&X\arrow[dl,"\alpha_1"']\arrow[d,"\alpha_0"]\arrow[dr,"\alpha_2"]\&\\
\underline{\mathrm{Ch}}(X)\arrow[dr,"p_1"] \& \mathcal C_{\mathbb N^o}\arrow[l,"\unsure{\pi|\Pi}_1"']\arrow[r,"\pi_2"]\arrow[d,"p_0"] \& \underline{\mathrm{Ch}}_\infty(X)\arrow[dl,"p_2"']\\
\&X\times X\&{.}
\end{tikzcd}}
Par 1) et $1')$ (concernant $\alpha_1$ et $\alpha_0$) et \unsure{le résultat|les résultats}
% archive sheet number: 418
